\documentclass[11pt]{scrartcl}
\usepackage[sexy]{evan}

\title{My First Difficult Olympiad Geometry Problem}
\author{Xinzhe Li}
\date{10 June 2026}

\begin{document}

\section*{My First Difficult Olympiad Geometry Problem}

\begin{mdframed}[style=mdpurplebox,frametitle={Problem statement}]
\textit{Source: We Love Geometry WeChat public account, Teacher Yao Jiabin.}

\medskip

Let $\triangle ABC$ be inscribed in a circle $\Gamma$.
Let $P$ be an interior point of $\triangle ABC$.

Let $E$ and $F$ be the feet of the perpendiculars from $P$ to $AC$ and $AB$,
respectively. Let $L$ be the intersection of line $EF$ with $\Gamma$ shown in
the figure.
Let line $AP$ meet $\Gamma$ again at $D$.
Finally, let $S$ be the foot of the perpendicular from $D$ to $BC$.

Prove that
\[
  \angle PLE = \angle SLD.
\]
\end{mdframed}

\begin{center}
\begin{asy}
size(13cm);

pen auxpen = rgb(0.45,0.45,0.45);

path localRightAngleMark(pair A, pair O, pair B, real scale=8) {
  real s = scale/40;
  pair u = unit(A-O);
  pair v = unit(B-O);
  return (O+s*u)--(O+s*(u+v))--(O+s*v);
}

real polarAngle(pair v) {
  return degrees(atan2(v.y, v.x));
}

path localSmallAngleArc(pair A, pair O, pair B, real r) {
  real a = polarAngle(A-O);
  real b = polarAngle(B-O);
  while (b-a > 180) {
    b -= 360;
  }
  while (b-a < -180) {
    b += 360;
  }
  return arc(O, r, a, b);
}

// ------------------------------------------------
// Basic configuration
// ------------------------------------------------
pair B = (0,0);
pair C = (6,0);
pair A = (5.4,4.7);
pair P = (3.55,1.65);

// Feet of perpendiculars
pair E = foot(P,A,C);
pair F = foot(P,A,B);

// Circumcircle Gamma
circle Gamma = circumcircle(A,B,C);
pair O = circumcenter(A,B,C);
real R = abs(A-O);

// L is the second intersection of EF with Gamma
pair[] X = intersectionpoints(line(E,F), Gamma);
pair L;
if (abs(X[0]-E) > abs(X[1]-E)) {
  L = X[0];
} else {
  L = X[1];
}

// D is the second intersection of AP with Gamma
pair[] Y = intersectionpoints(line(A,P), Gamma);
pair D;
if (abs(Y[0]-A) > abs(Y[1]-A)) {
  D = Y[0];
} else {
  D = Y[1];
}

// S is the foot from D to BC
pair S = foot(D,B,C);

// Simson line feet and auxiliary intersection
pair G = foot(D,A,B);
pair H = foot(D,A,C);
pair I = extension(A,L,G,H);

// ------------------------------------------------
// Evan-style drawing
// ------------------------------------------------

// Main circle: Evan-style filled circle
filldraw(Gamma, opacity(0.1)+lightcyan, blue);

// Main triangle
draw(A--B--C--cycle, linewidth(1.1));
draw(A--H, gray+dashed);

// Given construction lines
draw(E--F--L, lightred);
draw(A--D, lightred);

// Perpendiculars
draw(P--E, auxpen);
draw(P--F, auxpen);
draw(D--S, auxpen);
draw(D--G, auxpen);
draw(D--H, auxpen);

// Simson line and auxiliary lines
draw(I--H, purple);
draw(A--I, lightred);
draw(I--D, purple);

// Lines involved in the target equality
draw(L--P, blue);
draw(L--S, blue);
draw(L--D, blue);

// Right angle marks
draw(localRightAngleMark(P,E,C,8), auxpen);
draw(localRightAngleMark(P,F,A,8), auxpen);
draw(localRightAngleMark(D,S,C,8), auxpen);
draw(localRightAngleMark(D,G,A,8), auxpen);
draw(localRightAngleMark(D,H,A,8), auxpen);

// Small angle marks
draw(localSmallAngleArc(P,L,E,0.34), deepcyan);
draw(localSmallAngleArc(S,L,D,0.46), deepgreen);

// Points and labels
dot("$A$", A, dir(A));
dot("$B$", B, dir(225));
dot("$C$", C, dir(315));
dot("$D$", D, dir(270));
dot("$E$", E, dir(0));
dot("$F$", F, dir(135));
dot("$G$", G, dir(135));
dot("$H$", H, dir(0));
dot("$I$", I, dir(45));
dot("$L$", L, dir(180));
dot("$P$", P, dir(330));
dot("$S$", S, dir(90));

// Circle label
label("$\Gamma$", O + 1.08*R*dir(35), blue);

\end{asy}
\end{center}

\begin{proof}
Let $G$ and $H$ be the feet of the perpendiculars from $D$ to $AB$ and $AC$,
respectively. Thus
\[
  DG \perp AB, \qquad DH \perp AC.
\]
Since $D$ lies on the circumcircle of $\triangle ABC$, by the Simson line
theorem, the three feet $G,S,H$ are collinear.
All ratios of collinear segments below are understood as directed ratios.

Let
\[
  I \coloneq AL \cap GH.
\]
We first claim that $LP \parallel ID$. Indeed, since
\[
  PF \perp AB,\quad DG \perp AB,
\]
we have $PF \parallel DG$. Hence
\[
  \triangle AFP \sim \triangle AGD,
\]
and therefore
\[
  \frac{AF}{AG}=\frac{AP}{AD}.
\]
Similarly, since
\[
  PE \perp AC,\quad DH \perp AC,
\]
we have $PE \parallel DH$. Hence
\[
  \triangle AEP \sim \triangle AHD,
\]
and thus
\[
  \frac{AE}{AH}=\frac{AP}{AD}.
\]
Consequently,
\[
  \frac{AF}{AG}=\frac{AE}{AH},
\]
which implies
\[
  EF \parallel GH.
\]

Since $L,E,F$ are collinear and $I,G,H$ are collinear, this gives
\[
  LE \parallel IH.
\]
Moreover, applying the intercept theorem in triangle $AIG$, with
$L\in AI$, $F\in AG$, and $LF\parallel IG$, gives
\[
  \frac{AL}{AI}=\frac{AF}{AG}.
\]
Therefore
\[
  \frac{AL}{AI}=\frac{AP}{AD}.
\]
Since $A,L,I$ are collinear and $A,P,D$ are collinear, it follows that
\[
  LP \parallel ID.
\]

Now
\[
  LE\parallel IH,\qquad LP\parallel ID,\qquad EP\parallel HD.
\]
Hence
\[
  \triangle ELP \sim \triangle HID.
\]
In particular,
\[
  \dang PLE=\dang DIH.
\]

It remains to prove that
\[
  \dang DIH=\dang DLS.
\]
Since $A,L,D,C$ are concyclic and $A,L,I$ are collinear, we have
\[
  \dang ILD=\dang ALD=\dang ACD.
\]
Also,
\[
  DS\perp BC,\qquad DH\perp AC,
\]
so
\[
  \angle DSC=\angle DHC=90^\circ.
\]
Thus $S,C,H,D$ are concyclic. Hence
\[
  \dang HSD=\dang HCD.
\]
Since $A,C,H$ are collinear, we have
\[
  \dang ACD=\dang HCD.
\]
Since $I,S,H$ are collinear, we have
\[
  \dang ISD=\dang HSD.
\]
Therefore
\[
  \dang ACD=\dang ISD.
\]
Together with $\dang ILD=\dang ACD$, this gives
\[
  \dang ILD=\dang ISD.
\]
Thus $I,L,D,S$ are concyclic.

It follows that
\[
  \dang DLS=\dang DIS.
\]
Since $I,S,H$ are collinear, we have
\[
  \dang DIS=\dang DIH.
\]
Combining this with $\dang PLE=\dang DIH$, we obtain
\[
  \dang PLE=\dang DLS.
\]
Therefore, in ordinary angles,
\[
  \angle PLE=\angle SLD.
\]
This proves the desired result.
\end{proof}

\end{document}
